% =========================
% Chapter 10
% =========================
\chapter{Conclusion}
\label{ch:conclusion}

This work set out to answer a specific question: can an assistive system provide
reliable, context-aware guidance to blind users while minimizing cognitive load, by
activating perception only for explicit tasks? Lumen answers it constructively. A
blind user with nothing but an ordinary smartphone and a web address can ask for an
object and be guided to it, have their hand steered onto it and released with an
automatic ``grasp it,'' or name a destination room and be walked there
door-by-door --- with the system silent before the task, disciplined during it, and
silent again after.

The contributions claimed in Chapter~\ref{ch:introduction} hold up under the
project's own evaluation. The \emph{task-scoped perception paradigm} is enforced
structurally rather than by convention: a five-state finite state machine governs
every session, every transition is an explicit event, and a task engine cannot
outlive its state --- cancellation, completion, disconnection, and error all funnel
through the same teardown. \emph{Goal-directed exploration} demonstrates that
semantic room recognition, door detection, and a compass-anchored $360^\circ$ scan
can substitute for the localization a phone camera cannot provide, completing
multi-room journeys with no map, no markers, and no infrastructure. The
\emph{door perception funnel} shows what safety-grade perception on commodity
models requires: the custom detector reaches mAP@50 of $0.946$, yet the system
never acts on it raw --- every spoken door survived geometric gates, semantic
verification by a second model, and temporal evidence accumulation, a design forged
failure-by-failure against real rooms. \emph{Zero-install delivery} proved not
merely convenient but architecturally clarifying: one WebSocket, one server
process, one URL. And the \emph{testable architecture} kept the whole decision
layer honest --- 281 automated tests replay complete navigation journeys, safety
invariants, and fifty voice-command variants in under a second, on every change.

Beyond the artifacts, the project yields a methodological lesson worth stating
plainly. In assistive perception, the adversary is not the benchmark but the
world: blank walls that score as doors at 0.9 confidence, lace curtains that
defeat every geometric test, real doors that read as refrigerators, compasses that
lie precisely next to the appliances that define a kitchen. A system whose words a
blind person will act on must therefore be engineered so that no single model's
confidence is ever trusted alone, so that evidence must be repeated and localized
before it is spoken, and so that a false positive is always treated as more
expensive than a false negative --- a missed door costs a rescan; an invented one
points a person at a wall. Equally, the voice channel is not presentation but
interface: utterances that overlap, queue up, or contradict each other are
functional defects, and the speech discipline that prevents them required as much
engineering as the perception it narrates.

The system's boundaries are stated honestly in Chapter~\ref{ch:limitations} ---
no localization or door memory yet, direction hints that reach only the current
fork, an evaluation with sighted testers --- and each maps onto the
staged plan of Chapter~\ref{ch:future_work}, with deeper hint handling,
exploration memory, and a formal blind-user study as the decisive
next steps. What stands at the end of this project is a working demonstration
that structured, task-scoped assistance --- respecting the cognitive load,
autonomy, and safety of blind users --- can be delivered today, on hardware
already in the user's pocket, and a tested foundation on which that assistance
can now grow.